% The Interval + Path Types (Ch 1-2)
% Lines 120-560 of main.tex
% This file is for agent editing — will be merged back

\chapter{The Interval, Faces, and Connections}
\label{ch:interval}

Cubical type theory (CCHM, 2017; Cohen et al., 2018) is a type theory
where equality is witnessed by \emph{paths} --- functions from an
abstract interval.  We develop the theory in full, specialized to
program semantics.

\section{The Abstract Interval $\interval$}

\begin{definition}[The Interval]
The \emph{abstract interval} $\interval$ is a type with:
\begin{enumerate}
  \item Two \emph{endpoints}: $\mathsf{i0}, \mathsf{i1} : \interval$.
  \item \emph{Connections}: binary operations $\land, \lor : \interval
    \times \interval \to \interval$ satisfying:
    \begin{align}
      r \land \mathsf{i0} &= \mathsf{i0} &
      r \land \mathsf{i1} &= r \\
      r \lor \mathsf{i0} &= r &
      r \lor \mathsf{i1} &= \mathsf{i1}
    \end{align}
  \item \emph{Reversal}: $\lnot : \interval \to \interval$ with
    $\lnot \mathsf{i0} = \mathsf{i1}$ and
    $\lnot \mathsf{i1} = \mathsf{i0}$.
  \item The triple $(\interval, \land, \lor, \lnot, \mathsf{i0},
    \mathsf{i1})$ forms a \emph{De Morgan algebra}.
\end{enumerate}
\end{definition}

\begin{remark}[Why Not the Unit Interval?]
The interval $\interval$ is \emph{not} the topological unit interval
$[0,1]$.  It is a formal, syntactic object --- a free De Morgan algebra
on one generator.  Its purpose is to parameterize \emph{deformations}
between terms.  In our setting, a path $p : \interval \to (A \to B)$
with $p(\mathsf{i0}) = f$ and $p(\mathsf{i1}) = g$ is a formal
deformation from program $f$ to program $g$.
\end{remark}

\section{Faces and Face Lattice}

A \emph{face} is a constraint on interval variables.

\begin{definition}[Face Formulas]
The \emph{face lattice} $\mathbb{F}$ is the free distributive lattice
generated by atoms $i = 0$ and $i = 1$ for each interval variable $i$.
A face formula $\varphi \in \mathbb{F}$ is a conjunction/disjunction of
such atoms:
\[
  \varphi ::= (i = 0) \mid (i = 1) \mid \varphi_1 \land \varphi_2
  \mid \varphi_1 \lor \varphi_2
\]
\end{definition}

\begin{example}[Face of a Square]
A 2-dimensional cube (square) has four edges (1-faces) and four vertices
(0-faces).  The edge $i = 0$ is the left side; $j = 1$ is the top.
The vertex $i = 0 \land j = 1$ is the top-left corner.
\end{example}

In program equivalence, faces correspond to \emph{restrictions} of
the input space.  The face $\mathsf{tag}(p) = \mathsf{Int}$ restricts
to the integer fiber.

\section{The Face Map and Degeneracy}

\begin{definition}[Face Maps]
For a cubical term $t$ depending on interval variable $i$, the
\emph{face maps} are:
\begin{align}
  \partial_i^0(t) &= t[i \mapsto \mathsf{i0}] \\
  \partial_i^1(t) &= t[i \mapsto \mathsf{i1}]
\end{align}
These give the ``boundary values'' of $t$ at each endpoint of $i$.
\end{definition}

\begin{definition}[Degeneracy Maps]
The \emph{degeneracy map} $\sigma_i$ adds a ``trivial'' interval
dimension:
\[
  \sigma_i(t) = t \quad \text{(where $t$ does not depend on $i$)}
\]
A degenerate cube is one that is constant in some dimension.
\end{definition}

Face and degeneracy maps satisfy the \emph{cubical identities}:
\begin{align}
  \partial_j^b \circ \partial_i^a &= \partial_i^a \circ \partial_{j-1}^b
    \quad \text{for $j > i$} \\
  \sigma_j \circ \sigma_i &= \sigma_i \circ \sigma_{j+1}
    \quad \text{for $j \geq i$} \\
  \partial_j^b \circ \sigma_i &= \begin{cases}
    \sigma_{i-1} \circ \partial_j^b & j < i \\
    \mathrm{id} & j = i \\
    \sigma_i \circ \partial_{j-1}^b & j > i
  \end{cases}
\end{align}

These identities ensure that the cubical structure is coherent ---
different ways of computing the same face give the same result.

\begin{proposition}[The Cube Category $\square$]
\label{prop:cube-category}
The \emph{cube category} $\square$ has objects $[n] = \interval^n$
for $n \geq 0$ and morphisms generated by:
\begin{itemize}
  \item Face maps $\delta_i^\varepsilon : [n-1] \to [n]$ for $1 \leq i \leq n$
    and $\varepsilon \in \{0, 1\}$, inserting $\mathsf{i}\varepsilon$ at
    position~$i$.
  \item Degeneracy maps $\sigma_i : [n+1] \to [n]$ for $1 \leq i \leq n+1$,
    projecting away dimension~$i$.
  \item Connection maps $\gamma_i^\varepsilon : [n+1] \to [n]$ for
    $1 \leq i \leq n$, applying $\land$ (for $\varepsilon = 0$) or $\lor$
    (for $\varepsilon = 1$) to coordinates~$i$ and~$i+1$.
\end{itemize}
Subject to the cubical identities and the connection identities:
\begin{align}
  \delta_i^0 \circ \gamma_i^0 &= \sigma_i &
  \delta_i^1 \circ \gamma_i^0 &= \mathrm{id} \\
  \delta_i^0 \circ \gamma_i^1 &= \mathrm{id} &
  \delta_i^1 \circ \gamma_i^1 &= \sigma_i
\end{align}
\end{proposition}

\begin{proof}
The identities follow from the De Morgan algebra axioms of $\interval$.
For instance, $\delta_i^0 \circ \gamma_i^0$ applied to a term $t(r_1,
\ldots, r_{n+1})$ gives $t(r_1, \ldots, r_i \land r_{i+1}, \ldots)[r_i
\mapsto \mathsf{i0}] = t(\ldots, \mathsf{i0} \land r_{i+1}, \ldots)
= t(\ldots, \mathsf{i0}, \ldots)$, which equals $\sigma_i(t)$.
The identity $\delta_i^1 \circ \gamma_i^0$ gives
$t(\ldots, \mathsf{i1} \land r_{i+1}, \ldots) = t(\ldots, r_{i+1}, \ldots)
= \mathrm{id}(t)$.
\end{proof}

\begin{definition}[Cubical Sets]
\label{def:cubical-set}
A \emph{cubical set} is a presheaf $X : \square^{\op} \to \Set$.
Concretely, it consists of:
\begin{itemize}
  \item For each $n \geq 0$, a set $X_n$ of \emph{$n$-cubes}.
  \item For each face map $\delta_i^\varepsilon$, a restriction
    $\partial_i^\varepsilon : X_n \to X_{n-1}$.
  \item For each degeneracy map $\sigma_i$, an extension
    $s_i : X_n \to X_{n+1}$.
  \item For each connection, a diagonal map, all satisfying the
    dual cubical identities.
\end{itemize}
A \emph{morphism of cubical sets} $f : X \to Y$ is a natural transformation.
The category of cubical sets is written $\widehat{\square} \coloneqq [\square^{\op}, \Set]$.
\end{definition}

\begin{example}[Programs as a Cubical Set]
\label{ex:programs-cubical}
Fix a function type $A \to B$.  Define the cubical set $\mathcal{P}_{A \to B}$ by:
\begin{itemize}
  \item $\mathcal{P}_0 = \{$programs $f : A \to B\}$ (0-cubes are programs).
  \item $\mathcal{P}_1 = \{$formal deformations $p : \interval \to (A \to B)
    \mid p(\mathsf{i0}) = f,\, p(\mathsf{i1}) = g\}$ (1-cubes are equivalence
    witnesses).
  \item $\mathcal{P}_2$ consists of homotopies between equivalence proofs
    (coherence data).
\end{itemize}
The face maps $\partial_1^0, \partial_1^1 : \mathcal{P}_1 \to \mathcal{P}_0$
extract the source and target programs.  The degeneracy $s_1 : \mathcal{P}_0
\to \mathcal{P}_1$ sends each program to its reflexivity path.
\end{example}

\begin{theorem}[The Yoneda Embedding for $\square$]
\label{thm:yoneda-cube}
The representable presheaf $\square[-, [n]]$ is the standard $n$-cube.
By the Yoneda lemma, for any cubical set~$X$:
\[
  \mathrm{Hom}_{\widehat{\square}}(\square[-, [n]], X) \cong X_n
\]
Thus the $n$-cubes of $X$ are exactly the maps from the standard $n$-cube into $X$.
\end{theorem}

\begin{proof}
This is the standard Yoneda lemma applied to the functor category
$[\square^{\op}, \Set]$.  A natural transformation
$\alpha : \square[-, [n]] \Rightarrow X$ is determined by
$\alpha_{[n]}(\mathrm{id}_{[n]}) \in X_n$, and conversely every element
of $X_n$ determines such a transformation.
\end{proof}

\begin{remark}[Connection to the Implementation]
In the implementation (\texttt{cctt/theory.py}), the interval does not appear
as an explicit object.  Instead, its role is played by the \emph{parameter space}
of the programs being compared.  A ``path'' between programs $f$ and $g$ is
constructed by finding a rewrite chain $f = t_0 \to t_1 \to \cdots \to t_k = g$
in the HIT $\mathsf{PyComp}$ (Chapter~\ref{ch:hits}).  Each rewrite step $t_i \to t_{i+1}$
is a 1-cube in the cubical set of computations, and the chain is a composite
path formed by Kan filling (Definition~\ref{def:composition}).
\end{remark}


\chapter{Path Types and Function Extensionality}
\label{ch:paths}

\section{Path Types}

\begin{definition}[Path Type]
For a type $A$ and terms $a, b : A$, the \emph{path type}
$\Path_A(a, b)$ is:
\[
  \Path_A(a, b) = \{p : \interval \to A \mid p(\mathsf{i0}) = a
    \text{ and } p(\mathsf{i1}) = b\}
\]
A term $p : \Path_A(a, b)$ is a \emph{path from $a$ to $b$ in $A$}.
\end{definition}

\begin{definition}[Reflexivity]
For any $a : A$, the \emph{constant path} $\mathsf{refl}_a :
\Path_A(a, a)$ is:
\[
  \mathsf{refl}_a = \lambda i.\, a
\]
\end{definition}

\begin{definition}[Symmetry]
For $p : \Path_A(a, b)$, the \emph{inverse path} $p^{-1} :
\Path_A(b, a)$ is:
\[
  p^{-1} = \lambda i.\, p(\lnot i)
\]
\end{definition}

\begin{definition}[Path Composition via Kan Filling]
\label{def:composition}
For $p : \Path_A(a, b)$ and $q : \Path_A(b, c)$, the
\emph{composite path} $p \cdot q : \Path_A(a, c)$ is constructed
by the Kan filling operation.

Consider the open box:
\begin{center}
\begin{tikzcd}
  a \arrow[r, "p"] \arrow[d, dashed, "p \cdot q"'] &
  b \arrow[d, "q"] \\
  {} & c
\end{tikzcd}
\end{center}

The Kan filler closes the box, producing the composite path from $a$
to $c$.

Formally:
\[
  (p \cdot q)(i) = \hfilling^i_{j}
    \begin{cases}
      (j = \mathsf{i0}) &\mapsto a \\
      (j = \mathsf{i1}) &\mapsto q(i) \\
      (i = \mathsf{i0}) &\mapsto p(j)
    \end{cases}
\]
\end{definition}

\begin{theorem}[Path Composition is Well-Defined]
The Kan filling operation produces a well-defined term of type
$\Path_A(a, c)$.  The boundary conditions are:
\begin{itemize}
  \item $(p \cdot q)(\mathsf{i0}) = a$ (by the $i = \mathsf{i0}$
    face, which gives $p(\mathsf{i0}) = a$).
  \item $(p \cdot q)(\mathsf{i1}) = c$ (by the $j = \mathsf{i1}$
    face, which gives $q(\mathsf{i1}) = c$).
\end{itemize}
\end{theorem}

\section{Function Extensionality}

\begin{theorem}[Function Extensionality --- funext]
\label{thm:funext}
For $f, g : A \to B$:
\[
  \Path_{A \to B}(f, g) \simeq \prod_{x : A} \Path_B(f(x), g(x))
\]
That is, a path between functions is equivalent to a family of
pointwise paths.
\end{theorem}

\begin{proof}
\emph{(Forward direction)} Given $p : \Path_{A \to B}(f, g)$ and
$x : A$, define $\mathsf{happly}(p, x) = \lambda i.\, p(i)(x)$.
This is a path from $f(x)$ to $g(x)$.

\emph{(Backward direction)} Given
$h : \prod_{x : A} \Path_B(f(x), g(x))$, define
$\mathsf{funext}(h) = \lambda i.\, \lambda x.\, h(x)(i)$.
This is a path from $f$ to $g$.

The two directions are mutually inverse (up to higher paths).
\end{proof}

\begin{corollary}[Program Equivalence as Pointwise Path Family]
\label{cor:prog-equiv}
Two programs $f, g : A \to B$ are equivalent iff for every input
$x : A$, there is a path $p_x : \Path_B(f(x), g(x))$.  By funext,
this family assembles into a single path $f \equiv g$.
\end{corollary}

This is the \emph{fundamental theorem} of our approach: to prove
program equivalence, we construct pointwise paths (one per input)
and assemble them via funext.

\begin{lemma}[Path Inversion is an Involution]
\label{lem:path-inv-invol}
For any path $p : \Path_A(a, b)$, we have $(p^{-1})^{-1} = p$.
\end{lemma}

\begin{proof}
$(p^{-1})^{-1} = \lambda i.\, p^{-1}(\lnot i) = \lambda i.\, p(\lnot(\lnot i))
= \lambda i.\, p(i) = p$, using the involution property $\lnot \lnot r = r$
of the De Morgan algebra $\interval$.
\end{proof}

\begin{theorem}[Groupoid Structure of Paths]
\label{thm:path-groupoid}
For any type $A$, the path type structure forms a groupoid:
\begin{enumerate}
  \item \emph{Identity}: $\mathsf{refl}_a : \Path_A(a, a)$ for all $a$.
  \item \emph{Inverse}: $p^{-1} : \Path_A(b, a)$ for $p : \Path_A(a, b)$.
  \item \emph{Composition}: $p \cdot q : \Path_A(a, c)$ for $p : \Path_A(a, b)$,
    $q : \Path_A(b, c)$.
  \item \emph{Left unit}: $\mathsf{refl}_a \cdot p \sim p$.
  \item \emph{Right unit}: $p \cdot \mathsf{refl}_b \sim p$.
  \item \emph{Left inverse}: $p^{-1} \cdot p \sim \mathsf{refl}_b$.
  \item \emph{Right inverse}: $p \cdot p^{-1} \sim \mathsf{refl}_a$.
  \item \emph{Associativity}: $(p \cdot q) \cdot r \sim p \cdot (q \cdot r)$.
\end{enumerate}
Here $\sim$ denotes the existence of a 2-path (homotopy).
\end{theorem}

\begin{proof}[Proof sketch]
Each identity is constructed via Kan filling operations on appropriate
open boxes.  For instance, the left unit law $\mathsf{refl}_a \cdot p \sim p$
is witnessed by the 2-path:
\[
  \alpha(i, j) = \hfilling^i_k
  \begin{cases}
    (k = \mathsf{i0}) &\mapsto a \\
    (k = \mathsf{i1}) &\mapsto p(i) \\
    (j = \mathsf{i0}) &\mapsto \mathsf{refl}_a(k) = a \\
    (j = \mathsf{i1}) &\mapsto p(i)
  \end{cases}
\]
The boundary at $j = \mathsf{i0}$ gives $\mathsf{refl}_a \cdot p$ and
at $j = \mathsf{i1}$ gives $p$.  Associativity requires a 3-dimensional
filling argument; see Cohen et al.~(2018, \S 6.1).
\end{proof}

\begin{example}[Groupoid of Program Rewrites]
\label{ex:rewrite-groupoid}
Consider the type $\mathsf{PyComp}$ (Chapter~\ref{ch:hits}).  The groupoid
structure on its path type has a direct computational reading:
\begin{itemize}
  \item \emph{Identity}: every program is equivalent to itself
    (no rewrite needed).
  \item \emph{Inverse}: if $f$ can be rewritten to $g$, then $g$ can be
    rewritten back to $f$ (every axiom application is reversible).
  \item \emph{Composition}: rewrite chains can be concatenated.
\end{itemize}
The path search algorithm (\texttt{cctt/path\_search.py}) constructs
explicit witnesses of this groupoid structure: each \texttt{PathStep}
records a single axiom application, and a \texttt{PathResult} is a
chain of such steps, which is exactly a composite path
$p_1 \cdot p_2 \cdots p_k$ in $\mathsf{PyComp}$.
\end{example}

\begin{definition}[Dependent Path Type]
\label{def:dep-path}
For a type family $P : A \to \Type$ and a path $p : \Path_A(a, b)$,
the \emph{dependent path type} (path-over) is:
\[
  \mathsf{PathP}(P, p, u, v) \coloneqq \{q : \prod_{i : \interval} P(p(i))
  \mid q(\mathsf{i0}) = u,\; q(\mathsf{i1}) = v\}
\]
where $u : P(a)$ and $v : P(b)$.
\end{definition}

\begin{proposition}[PathP Reduces to Path When $P$ is Constant]
If $P \equiv \lambda a.\, B$ is a constant family, then
$\mathsf{PathP}(P, p, u, v) \simeq \Path_B(u, v)$.
\end{proposition}

\begin{proof}
When $P$ is constant, $P(p(i)) = B$ for all $i$, so a dependent path
$q : \prod_{i : \interval} B$ with endpoints $u, v$ is exactly an
ordinary path in~$B$.
\end{proof}

\section{Transport and Substitution}

\begin{definition}[Transport]
Given a type family $B : A \to \Type$ and a path $p : \Path_A(a, a')$,
\emph{transport} along $p$ is:
\[
  \transp(B, p) : B(a) \to B(a')
\]
defined by $\transp(B, p)(b) = \comp^i(B(p(i)), [\,], b)$.
\end{definition}

Transport is essential for \emph{type-changing} equivalences.  When
a program returns $\mathsf{int}$ in one version and $\mathsf{bool}$
in another, transport along the coercion path $\mathsf{int} \to
\mathsf{bool}$ (via Python's truthiness) relates the two outputs.

\section{Higher Paths and Truncation}

\begin{definition}[Higher Paths]
A \emph{2-path} (homotopy) between paths $p, q : \Path_A(a, b)$ is:
\[
  \alpha : \Path_{\Path_A(a,b)}(p, q)
\]
This is a ``path between paths'' --- a square
$\alpha : \interval \times \interval \to A$ with boundary:
$\alpha(\mathsf{i0}, -) = p$, $\alpha(\mathsf{i1}, -) = q$,
$\alpha(-, \mathsf{i0}) = \mathsf{refl}_a$,
$\alpha(-, \mathsf{i1}) = \mathsf{refl}_b$.
\end{definition}

\begin{definition}[Propositional Truncation]
A type $A$ is \emph{propositional} (or a \emph{mere proposition}) if
any two elements are connected by a path:
\[
  \mathsf{isProp}(A) = \prod_{a, b : A} \Path_A(a, b)
\]
The \emph{propositional truncation} $\|A\|$ makes any type
propositional by adding a path between every pair of elements.
\end{definition}

\begin{remark}[Truncation for Equivalence]
Program equivalence is a \emph{proposition}: either $f \equiv g$ or
$f \not\equiv g$, and there is at most one proof (up to higher paths).
We work in the propositionally truncated path type
$\|\Path_{A \to B}(f, g)\|$.  This means we don't distinguish between
different proofs of equivalence --- we only care about the
\emph{existence} of a path.
\end{remark}


